\section{课程不是名人目录：先建立读课方法}

这节开场课没有从 GPU 型号、分布式训练或云计算 API 开始，而是先回答一个更基础的问题：面对一门由大量产业领袖参与的课程，学生应当怎样学习。讲者把自己称为整季演出的 opening act，把后续嘉宾称为 headliners，又主动接住 ``AI Coachella'' 的调侃。这个开场在提醒读者，阵容本身不是知识；真正有价值的是比较不同角色怎样定义问题、怎样暴露约束，以及他们的利益位置会怎样影响判断。

\begin{figure}[H]
\centering
\includegraphics[width=0.9\textwidth]{images/00-02-45-life-scaling-laws.jpg}
\caption{讲者借用 scaling laws 的语言，给出关于关系、乐趣与长期影响力的生活启发式。}
\figsource{Stanford Online 官方视频 00:02:45，`images/00-02-45-life-scaling-laws.jpg`。}
\end{figure}

\begin{knowledgebox}{读图：为什么技术课先讲 ``life scaling laws''}
这页并不是把人生伪装成可精确预测的训练曲线。讲者强调自己的经验法则来自个人观察：认真投入工作，同时保留乐趣，并和愿意长期相处的人建立关系。关键词是 \term{empirical}，也就是经验上有效，但不保证对每个人、每个阶段都成立。把这层限制说清楚，反而为后面讨论 AI scaling 做了铺垫：曲线可以描述已经观察到的规律，不能自动证明未来会无限延伸。
\end{knowledgebox}

\begin{warningbox}{不要把比喻当成定律}
``人生像训练任务'' 只能帮助整理决策变量，不能给人生成一个可优化的单一 objective。健康、关系、责任与影响力之间没有统一标尺。讲者保留 ``have fun with people you enjoy'' 这条朴素经验，正是为了避免把所有选择压缩成履历、薪酬或影响力数字。
\end{warningbox}

\subsection{Preparedness：课程要训练什么}

本节接着把学习目标从“听到名人观点”改成 \term{preparedness}，即面对真实系统、组织与社会约束时的准备度。准备度不等于记住嘉宾的结论，更接近三项能力：能把一个宏大主张拆成可检查的技术与经济假设；能看见自己所在层级之外的依赖；能在证据不足时保留不确定性。后续十周的嘉宾课因此应当被当作一组立场不同的现场案例，而不是一套统一答案。

\begin{figure}[H]
\centering
\includegraphics[width=0.86\textwidth]{images/00-07-45-introduction-to-anj.jpg}
\caption{Anjney Midha 的学习与职业路径，以及他在多家 AI 机构早期阶段的参与经历。}
\figsource{Stanford Online 官方视频 00:07:45，`images/00-07-45-introduction-to-anj.jpg`。}
\end{figure}

这页应当作为观点来源说明来读。讲者的经历覆盖 Stanford 数学、计算科学、经济学与生物信息学训练，也涉及 Discord、Anthropic、Mistral、Black Forest Labs、AMP 与 Periodic Labs 等组织。这样的履历让他更敏感于研究成果如何变成公司、产品和基础设施，同时也意味着他的判断来自投资、创业与部署视角。读者需要吸收实践经验，也要保留对样本偏差和利益位置的审计。

\begin{knowledgebox}{课堂提示：课程目标不是实习优化}
讲者明确区分 internship 与 preparedness。前者容易把每一讲压成“这家公司招什么人”；后者要求理解一家组织为什么在某个层级投入、它受哪些上游资源约束、又把哪些成本转嫁给下游。做课程项目时，问题定义、证据链与可验证原型因此比追逐热门名词更重要。
\end{knowledgebox}

\begin{importantbox}{本讲的阅读框架}
后文所有案例都用五个问题审计：谁掌握稀缺资源；资源如何进入产品闭环；扩张依赖什么反馈；当前证据能够证明什么；标准、市场或制度由谁协调。这个框架既适用于模型训练，也适用于 GPU 市场、云服务和公共基础设施。
\end{importantbox}

\subsection{本章小结}

开场中的生活经验、讲者背景和嘉宾阵容共同规定了课程方法：重视一手实践，但不崇拜身份；使用经验规律，但不把它们伪装成确定预测；把学习目标从短期机会扩大到跨层系统判断。下面进入本讲真正的技术主线：AI 为什么正在重新打开一套曾经相对稳定的系统栈。

